% =============================================================================
% 2. LITERATURE REVIEW
% =============================================================================
\section{Literature Review}
\label{sec:literature}

This section examines four bodies of literature that converge in the present
study: (i) the Kirkpatrick model and its modern critiques; (ii) aviation
maintenance training under regulatory pressure; (iii) human factors and
socio-technical perspectives on maintenance error; and (iv) technological and
economic enablers for next-generation evaluation. Section~\ref{subsec:gap}
synthesises the gap.

\subsection{Kirkpatrick's model: theoretical evolution and modern critiques}
\label{subsec:kirkpatrick}

Originally introduced in 1959 and consolidated in subsequent revisions, the
four-level model of training evaluation \citep{Kirkpatrick2010} remains the
most cited framework worldwide for assessing learning interventions. The
levels are widely understood as
\textit{Reaction} (participant satisfaction), \textit{Learning} (knowledge or
skill acquisition), \textit{Behaviour} (transfer of training to the workplace)
and \textit{Results} (organisational outcomes). The simplicity of the model
explains its longevity, but it has also drawn substantial criticism for
treating the four levels as a linear causal chain
\citep{Patel2023,Fernandez2024}.

Contemporary work has refined the model along three axes. First, attribution
methods such as SHAP-based causal analysis and quasi-experimental design have
been proposed to isolate the unique contribution of training when many other
variables change simultaneously \citep{Fernandez2024}. Second, predictive
extensions have been developed to forecast Level~3 and Level~4 outcomes from
Level~1 and Level~2 signals, enabling more proactive management
\citep{Peterson2025}. Third, integrative variants have linked Kirkpatrick to
safety culture, human-capital metrics and sustainability indicators
\citep{Sousa2025,Yang2024}. Despite these advances, empirical applications in
regulated heavy-asset sectors remain limited.

\subsection{Aviation maintenance training: regulatory imperatives and empirical shortfalls}
\label{subsec:aviation-training}

Aviation \MRO{} is regulated by a multi-level architecture that originates in
the \ICAO{} Convention on International Civil Aviation, is implemented by the
\FAA{}, \EASA{}, \ANAC{} and other civil aviation authorities (\CAA{}), and
operationalised by individual airlines, airports and \MRO{} facilities
through their Safety Management System (\SMS{}) manuals
\citep{FAA2024,EASA2025}. In parallel, the aerospace quality standard
\AS{} imposes specific requirements on the training process under its
clauses on resources (7.1), competence (7.2), awareness (7.3) and
communication (7.4) \citep{Oschman2017}.

Empirical work consistently reports that compliance with these requirements
does not automatically translate into reduced maintenance error.
\citet{Nakamura2024} document a \emph{compliance paradox} whereby
organisations satisfy regulatory documentation requirements while their
underlying behavioural and outcome metrics remain stagnant.
\citet{ORconnor2024} highlight the moderating role of safety culture in
determining whether training transfers to operational behaviour, while
\citet{Boyd2015} show that maintenance-related accident rates have plateaued
over the past two decades despite expanded training budgets. Recent
reviews specific to \MRO{} economics indicate that quality-cost reductions
from training are achievable but require systematic evaluation beyond
Level~1 \citep{Zhang2025,Garcia2024}.

\subsection{Human factors and socio-technical systems in maintenance}
\label{subsec:human-factors}

The contemporary understanding of maintenance error rests on the recognition
that errors are properties of socio-technical systems rather than properties
of individuals. \citet{Reason2016} provides the canonical model of latent and
active failures, recently extended by \citet{Leveson2024} into systems-theory
formulations that emphasise hierarchical control and feedback loops.
\citet{Dismukes2023} adds cognitive depth, showing that procedural deviations
in maintenance frequently originate from attentional bottlenecks rather than
from lack of knowledge, which has direct implications for the design of
training and assessment.

The relevance of this stream of literature for the present case study is
twofold. First, it justifies why a Level~2 (Learning) outcome alone does not
guarantee a Level~3 (Behaviour) outcome --- attention, fatigue, time pressure
and supervision practice all intervene. Second, it grounds the choice of
combining post-training questionnaires with direct on-the-job observation in
the data collection design (Section~\ref{sec:methodology}).

\subsection{Technological enablers for next-generation evaluation}
\label{subsec:tech-enablers}

Recent work indicates that digital tools --- notably digital twins for
procedural rehearsal \citep{Garcia2024} and causal-attribution methods
\citep{Fernandez2024} --- can extend training evaluation beyond
paper-based satisfaction surveys. While these enablers are not used in
the present case study, they are revisited in
Section~\ref{subsec:limitations} as future-research directions for
disentangling training effects from contemporaneous confounders and for
continuous Level~3 measurement.

\subsection{Bridging theory to practice: the safety--sustainability nexus}
\label{subsec:bridge}

The most recent literature explicitly links human-capital development in
high-reliability organisations with cleaner-production and circular-economy
outcomes \citep{Sousa2025,Yang2024,Chen2023}. The argument runs as follows:
every prevented error reduces the consumption of raw material, energy and
labour associated with rework, scrap, warranty and concession. When the
prevention mechanism is training, the resulting savings should be visible in
quality-cost accounts and in \ESG{} disclosures.

\subsection{Research gap}
\label{subsec:gap}

Three observations summarise the gap targeted by this paper:

\begin{enumerate}[label=(\roman*)]
  \item Despite the maturity of Kirkpatrick, empirical operationalisation
        beyond Level~2 in \AS{}-certified \MRO{} facilities is rare and
        the few existing studies report only safety or only economic
        outcomes, not both.
  \item Regulatory frameworks (\FAA{}, \EASA{}, \ANAC{}) mandate training but
        do not prescribe how to evaluate it across the four Kirkpatrick
        levels, creating a wide variance in practice.
  \item The literature linking training effectiveness to sustainability
        outcomes in aviation \MRO{} is incipient: the explicit translation of
        Levels~3 and~4 into material, energy and waste indicators has not yet
        been documented in a single longitudinal case.
\end{enumerate}

The present case study addresses these three observations directly through
the design described in Section~\ref{sec:methodology}.
